\section{Related Work}

This study lies at the intersection of four lines of work: noticed-case retrieval, THUIR's structurally informed legal pretraining, retrieval interaction architectures and ranking objectives, and long-context encoders with dynamic negative design.

\subsection{Legal Case Retrieval and THUIR@COLIEE 2023}

COLIEE Task~1 is not ordinary similarity search. The system must identify the prior cases that truly support the reasoning of the query judgment. THUIR@COLIEE 2023 \cite{li2023thuir} built a strong system around SAILER, reference sentence extraction, post-processing, and learning to rank. Its most important contribution is a structure-aware legal pretraining method: the judgment is divided into Fact, Reasoning, and Decision; a deeper encoder reads Fact and produces a `[CLS]` embedding; that embedding is then injected into two shallower decoders that reconstruct masked Reasoning and Decision. The overall objective is the sum of three masked language modeling losses, with lighter masking on Fact and heavier masking on Reasoning and Decision, so that the encoder is encouraged to learn a strong case representation \cite{li2023thuir}. Our work follows THUIR's motivation, but asks two more direct questions: whether citation-context queries are still suitable once long-context encoders are available, and whether dense retrieval is mainly bottlenecked by truncation.

\subsection{Retrieval Architectures and Ranking Objectives}

Among neural retrieval architectures, dual encoders encode the query and the candidate separately and compare them with vector similarity, making offline indexing possible and keeping first-stage retrieval efficient \cite{karpukhin2020dpr}. Cross encoders concatenate the query and the candidate into the same Transformer, allowing full token-level interaction but requiring separate inference for every query--document pair, which is much more expensive \cite{nogueira2019bertpr}. Late-interaction models such as ColBERT sit between the two: they encode query and document tokens separately and use cheaper token-level matching at scoring time \cite{khattab2020colbert}.

This distinction matters directly in our setting. Because both the query and the candidate are long judgments, a cross encoder would quickly exceed the token budget or incur prohibitive memory cost on a single RTX 4080 SUPER 16GB. We therefore use a dual encoder for first-stage dense retrieval and leave cross encoders and late interaction to future work. On the ranking side, pointwise methods treat each query--document pair as an independent example, pairwise methods learn which document should be ranked above another, and listwise methods optimize over the whole ranked list \cite{burges2005ranknet,cao2007listnet,liu2009ltr}. Our dual-encoder contrastive training is closer to listwise learning: although similarities are computed pairwise between the query and each candidate, the InfoNCE loss jointly normalizes one positive against 15 negatives in the same softmax, so the optimization target is the relative ordering of a whole local candidate set rather than isolated binary pairs. By contrast, LightGBM LambdaRank is pairwise in its update form because it builds gradients from document swaps, but those pairwise lambdas are weighted by their effect on overall NDCG, so its objective remains strongly listwise in spirit.

\subsection{Long-Context Encoders and ModernBERT}

ModernBERT \cite{warner2024modernbert} is a recent encoder-only Transformer designed to improve context length, memory efficiency, and inference speed at the same time. Its contribution is not merely a larger context window. Architecturally, it uses RoPE, pre-norm, and GeGLU while removing most bias terms; for long-context efficiency, it alternates global attention with local sliding-window attention, using global attention every third layer and 128-token local attention elsewhere; and at the implementation level it combines unpadding, Flash Attention 2/3, and \texttt{torch.compile} to accelerate long-context training and inference \cite{warner2024modernbert}. For this study, ModernBERT is valuable because it provides a practical testbed for the hypothesis that truncation is a central bottleneck in legal dense retrieval.

\subsection{Dynamic Negatives and Legally Plausible Candidate Scope}

Contrastive dense retrieval depends heavily on negative quality. If negatives are too easy, the model learns little; if they are fixed and no longer aligned with the current model state, training can also fail. In our early experiments, a THUIR-style BM25 top-100 hard-negative pool with random sampling of 15 negatives per update almost never converged. This suggests that in legal case retrieval, negative difficulty cannot be approximated by lexical rank alone.

Legal retrieval also has a second issue that is less prominent in general web search: without temporal scope constraints, semantically related but future judgments may be mixed into the negative pool. Such cases are not labeled as noticed cases, but that does not mean they are truly irrelevant; the judge may simply have been unable to cite a case from the future. We therefore combine dynamic negative sampling with query-specific year scope constraints so that the model learns negatives only from legally plausible candidates.
